\section{Variáveis Aleatórias}

Variáveis aleatórias (v.a) são representações de características numéricas de
experimentos aleatórios. Essa visão extende-se desde dos casos mais simples, a exemplo de
quando o próprio resultado é o atributo numérico de interesse, quanto a características
mais complexas, como a soma dos lados de dados, quantas vezes lançamos uma moeda até
obtermos duas caras ou o número de peças defeituosas produzidas após um dia de trabalho.
As v.a's podem ser formalizadas como funções que associam cada resultado do
espaço amostral a um número real.

\begin{def:va}
Se $\Omega$ é um espaço amostral, então uma variável aleatória $X$ é uma aplicação
de $\Omega$ em $\mathbb{R}$.
\end{def:va}

% É razoável presumir que a probabilidade de uma variável assumir uma realização específica
% é igual à probabilidade do evento que reune os resultados que levam a variável a ter
% essa realização. Por exemplo, se o experimento consiste em lançar dois
% dados, e $X$ é a soma dos dados, então a probabilidade da variável ter realização 7, que
% denotaremos por $X = 7$, deve ser igual à probabilidade do evento dos resultados em que a soma
% dos lados é 7, ou seja, do evento $\{(1, 6), (2, 5), (4, 3), (5, 2), (6, 1)\}$.

Introduzamos uma forma que nos permite definir eventos a partir de variáveis aleatórias
e os valores que elas mapeiam um determinado resultado no espaço amostral.

\begin{def:evento va}
Seja $X: \Omega \to \mathbb{R}$ uma variável aleatória. Se $A \subset \mathbb{R}$, então
o evento denotado por $\{X \in A\}$ é equivalente a
\[
	\{\omega \in \Omega \mid X(\omega) \in A\}
\]
\end{def:evento va}

Um caso especial dessa forma de definir eventos por meio de variáveis são as realizações,
que representam os casos em que queremos dizer qual foi o valor ao qual a variável mapeou
o resultado do experimento aleatório.

\begin{def:realizacao}
Seja $X: \Omega \to \mathbb{R}$ uma variável aleatória. Chamamos de realização de
$X$ igual a $x$ o evento definido por $\{\omega \in \Omega \mid X(\omega) = x\}$.
\end{def:realizacao}

Muitos outros predicados podem ser úteis para sintetizar eventos em termos dos valores
assumidos por uma variável. Por exemplo, para dizer que uma variável $X$ teve realização
menor do que $x$, então representamos por $\{X < x\}$. Ou que a realização de $X$ está
no intervalo $[a, b]$, denotamos por $\{X \in [a, b]\}$.

O conjunto dos valores cuja probabilidade da realização é não nula é chamado de suporte.
Por exemplo, no lançamento de um dado, as realizações possíveis são os inteiros de 1
a 6, enquanto os demais inteiros tem probabilidade nula. Coletando dados sobre o peso
corporal de uma população, o suporte é, em tese, qualquer valor positivo.

\begin{def:suporte}
Seja $X: \Omega \to \mathbb{R}$ uma variável aleatória.
O suporte de $X$ é o conjunto
\[
	\Omega_X = \{x \in \mathbb{R} \mid \forall U \text{ vizinhança aberta de } x, P(X \in U) > 0 \}
\]

\end{def:suporte}
% As variáveis tem associadas a si uma \textbf{função de distribuição de probabilidade}.
% Em síntese, uma função de distribuição de probabilidade nos dá uma ideia do quão
% frequente é um valor ou intervalos de valores de uma variável aleatória, além de poder
% informar sobre valores centrais e variabilidade da variável. A interpretação bayesiana
% para essas funções seria como a nossa crença modela a confiança \textit{a priori} em cada
% resultado -- ou intervalo deles -- possível do experimento.

% \begin{def:funcao de probabilidade}
% Seja $X: \Omega \to \mathbb{R}$ e $f: \mathbb{R} \to \mathbb{R}$. Dizemos que $f$ é uma
% função de probabilidade para
% \end{def:funcao de probabilidade}


% Uma família de funções de
% probabilidades, cujos membros diferenciam-se a menos do valor de um ou mais parâmetros,
% são chamadas de \textbf{distribuições de probabilidade}.

% \begin{def:suporte discreta}
% Seja $X: \Omega \to \mathbb{R}$ uma variável aleatória e $f: \mathbb{R} \to \mathbb{R}$ sua
% função de distribuição de probabilidade. O suporte de $X$, denotado por $\Omega_X$, é dado
% por
% \[
% 	\Omega_X = \{x \in \mathbb{R} \mid f(x) > 0 \}
% \]
% \end{def:suporte discreta}
% Portanto, o suporte de uma variável aleatória é formada pelos valores que são realizações
% possíveis dessa variável. Por exemplo, no lançamento de um dado de 6 faces, o suporte
% é formado pelos inteiros de 1 a 6; o número de pessoas com diploma numa população não
% pode ser negativo.
